\section{NetScope --- DNS Anomaly Detection}

\subsection*{Introduction}
NetScope addresses a specific cybersecurity problem: telling an ordinary DNS
lookup apart from one that statistically resembles data exfiltration or
command-and-control traffic. DNS tunneling encodes data --- credentials,
files, stolen records --- into the domain name of a DNS query, then sends
it to a server the attacker controls; because DNS traffic is high-volume
and rarely inspected in depth, this is hard to catch by hand, and hard to
catch with a static blocklist since the domain name itself may never have
been seen before. Unlike DarkWatch and CrisisBot, which need no trained
model at all, NetScope is a genuinely trained supervised classifier ---
one of several agents in \ProjectName{} built around a model learned from
labeled data rather than deterministic rules or an external lookup. What
makes it architecturally interesting is that it is really two things
joined together: an \textbf{offline training process} (a notebook, run on
Kaggle against a public labeled dataset, whose only job is to produce a
model file) and an \textbf{always-on production service}
(\texttt{netscope\_project/}, in this repository) that loads that model and
applies it continuously to live-looking traffic. Every technical claim
about the production code in this section was verified directly against
\texttt{src/netscope\_agent/}; claims that only the notebook can settle are
flagged explicitly where the notebook itself is not present in this
repository to re-check.

\subsection*{Role in the \ProjectName{} Architecture}
NetScope is an independent web service (FastAPI, its own port) launched
alongside the other four agents. Like NetScope's siblings, it does not
communicate directly with any other detection agent; its only outbound link
is to CrisisBot, triggered automatically once a score clears a fixed bar.
Three separate entry points can feed a DNS query into the same pipeline,
which mirrors the pattern already seen with DarkWatch's own three entry
points:

\begin{itemize}[nosep]
\item \textbf{\texttt{POST /analyze}} --- on-demand scoring, features already supplied by the caller (used by the dashboard).
\item \textbf{\texttt{POST /webhook/traffic-event}} --- a push endpoint for an external firewall or proxy, protected by a shared secret checked with \texttt{hmac.compare\_digest}. If the caller only supplies a domain name, the 16 features are computed live instead of being rejected.
\item \textbf{\texttt{auto\_scan.py}} --- a pull-based autonomous loop that samples a domain roughly every 25 seconds and can act on its own once a score crosses a threshold, the same principle as DarkWatch's \texttt{agent.py}.
\end{itemize}

Whichever entry point fires, the query passes through the same four stages
--- perception, feature extraction, scoring, decision --- before reaching
one of two outcomes, shown in Figure~\ref{fig:netscope-pipeline}.

\begin{figure}[htbp]
\centering
\begin{tikzpicture}[
  box/.style={draw, rounded corners, align=center, minimum width=4.6cm, minimum height=1.05cm, font=\small, fill=white},
  accent/.style={box, draw=blue!40!black, line width=1pt},
  arr/.style={-{Latex[length=2mm]}, thick},
  node distance=7mm and 8mm
]
\node[box] (entries) {3 entry points\\[1pt]\footnotesize \texttt{/analyze}, \texttt{/webhook/traffic-event}, \texttt{auto\_scan.py}};
\node[box, below=of entries] (perception) {Perception\\[1pt]\footnotesize raw domain name};
\node[box, below=of perception] (features) {Feature extraction\\[1pt]\footnotesize 16 features, live if needed (\texttt{live\_features.py})};
\node[box, below=of features] (scoring) {Scoring\\[1pt]\footnotesize LightGBM $\rightarrow$ \texttt{ml\_score} (0--100)};
\node[box, below=of scoring] (decision) {Decision\\[1pt]\footnotesize thresholds 50 / 80};
\node[box, below left=10mm and 6mm of decision] (noaction) {No automatic action\\[1pt]\footnotesize $<$50 normal, 50--80 suspect (dashboard)};
\node[accent, below right=10mm and 6mm of decision] (highscore) {score $\geq 80$};
\node[accent, below=of highscore] (crisisbot) {CrisisBot\\[1pt]\footnotesize unauthenticated webhook};
\node[accent, below=of crisisbot] (incident) {Incident response\\[1pt]\footnotesize CrisisBot's 5-step runbook};
\node[box, below=38mm of decision] (dash) {NetScope dashboard};

\draw[arr] (entries) -- (perception);
\draw[arr] (perception) -- (features);
\draw[arr] (features) -- (scoring);
\draw[arr] (scoring) -- (decision);
\draw[arr] (decision.south) -| (noaction.north);
\draw[arr] (decision.south) -| (highscore.north);
\draw[arr] (highscore) -- (crisisbot);
\draw[arr] (crisisbot) -- (incident);
\draw[arr] (noaction.south) |- (dash.west);
\draw[arr] (incident.south) |- (dash.east);
\end{tikzpicture}
\caption{NetScope --- pipeline (\texttt{src/netscope\_agent/}), styled to mirror DarkWatch's Figure~\ref{fig:darkwatch-pipeline-main}: a single decision point branches into a plain path (dashboard only) and an accented path that reaches CrisisBot automatically once the score clears 80.}
\label{fig:netscope-pipeline}
\end{figure}

Worked example: an external firewall sends
\texttt{POST /webhook/traffic-event?x\_webhook\_secret=...} with
\texttt{\{"domain": "xj4k9z-relay.info"\}} and no features attached. The
shared secret is checked first; since no \texttt{features} dict was
supplied, \texttt{extract\_all\_features\_live()} computes all 16 (10 from
the domain string alone, 6 from a real DNS resolution). Say the model
returns an \texttt{ml\_score} of 87: it clears both the 50 verdict
threshold (``suspect'') and the 80 auto-action threshold, so the host is
isolated automatically, with no human in the loop for this particular case.

Unlike DarkWatch (no trained model at all) or PhishHunter (an LLM in the
decision loop), NetScope is a classic supervised classification pipeline at
its core: the model is trained offline, exported, and loaded once at
service startup to score requests in real time. The LLM used later in this
section never touches the classification decision --- it only writes about
a result the model has already produced.

\subsection*{Technical Approach --- Two Halves of One Pipeline}
As in the introduction, NetScope splits cleanly into an offline half and an
online half, and the codebase implements this as a genuine handoff rather
than a blurred line:

\begin{itemize}[nosep]
\item \textbf{Offline training} (a Kaggle notebook, not stored in this repository): loads the public dataset, selects and justifies 16 features out of 44 raw columns, compares five model families under cross-validation, tunes the winner with Optuna, and exports \texttt{netscope\_model.pkl} and \texttt{netscope\_features.pkl}.
\item \textbf{Online serving} (\texttt{src/netscope\_agent/scorer.py} and the modules around it): loads those two exported files once at startup, reconstructs the same 16 features from nothing but a domain name when a real caller does not supply them, and applies the decision, explanation, and response logic on every query.
\end{itemize}

The rest of this section follows that same order --- dataset, features,
model, then the production logic built on top of it --- rather than
treating the notebook and the running service as two separate stories.

\subsection*{Dataset and Data Preparation}
The model was trained on the \textbf{CIC-Bell-DNS-Exfiltration 2021}
dataset (\texttt{dhoogla/cicbelldnsexf2021} on Kaggle), a public dataset
built specifically for DNS exfiltration detection. It ships in six variants
--- two cleaning methods (``minimalistic'' and ``opinionated'') crossed with
three volumes (``clean-Heavy'', ``clean-BenignExtra'', ``clean-Light''). The
training notebook loads the \texttt{minimalistic-clean-Light.parquet}
variant; only its two exported artifacts are copied into \texttt{models/}
for the production service to load (see \texttt{models/README.md}), so this
section reports what the notebook's own cells and markdown commentary
state rather than something re-run for this audit.

The raw file holds \textbf{102,762 DNS queries} described by \textbf{44
columns}, including per-record-type frequency statistics
(\texttt{A\_frequency}, \texttt{NS\_frequency}, \texttt{CNAME\_frequency},
and others), domain name structure statistics, TTL statistics, and two
original label columns: \texttt{GlobalClass} and \texttt{SubClass}.
\texttt{GlobalClass} takes the values \texttt{"Light Benign"} and
\texttt{"Light Attack"}. The notebook reduces this to a binary
\texttt{target} column, with the reasoning spelled out directly in a
markdown cell: NetScope needs to answer one simple question --- is this
domain suspicious or not --- so \texttt{GlobalClass} is collapsed into two
values (0 = Benign, 1 = Attack). Across the 102,762 rows, the real split is
\textbf{60,088 Benign} against \textbf{42,674 Attack}, a moderate imbalance
(58.5\,\% / 41.5\,\%) rather than an extreme one. This binary
\texttt{target} is exactly what \texttt{scorer.py} reproduces in
production --- \texttt{model.predict\_proba(x)[0][1]} is the model's
estimate of \texttt{target = 1} (Attack) for one query, turned into the
0--100 \texttt{ml\_score} every downstream stage relies on.

Before training, the notebook checks for missing values
(\texttt{isnull().sum()}) and duplicate rows (\texttt{duplicated().sum()}).
No missing values were found on the label-related columns
(\texttt{GlobalClass}, \texttt{SubClass}, \texttt{target}); the exact
duplicate-row count the notebook prints could not be independently
re-verified in this session, since the notebook itself is not stored in
this repository and only runs on Kaggle.

\subsection*{DNS Features}
Of the 44 raw columns, 16 were kept for training. The notebook's own
justification: keep the numeric columns that describe a domain name and its
DNS query statistically --- name entropy, length, subdomain count, share of
digits and special characters, TTL statistics --- rather than the raw
domain string itself, which a model like this cannot use directly. These
exact same 16 names, in this exact order, are what
\texttt{netscope\_features.pkl} exports and what \texttt{NetScopeScorer}
demands from every caller in production (\texttt{score\_domain()} raises an
explicit error, rather than silently filling a NaN, if any one of the 16 is
missing) --- the feature contract is not renegotiated between training and
serving.

\begin{table}[h]
\centering
\small
\begin{tabular}{@{}p{3cm}p{4.8cm}p{5.5cm}@{}}
\toprule
\textbf{Feature} & \textbf{What it measures} & \textbf{Why it matters for DNS security} \\
\midrule
\texttt{entropy} &
Shannon entropy of the full domain name &
A name generated to encode data reads as more random than one a human would
choose. \\

\texttt{rr\_name\_entropy} &
Entropy of the name inside the DNS resource record &
Same signal as \texttt{entropy}, computed one layer closer to the actual DNS
packet. \\

\texttt{rr\_name\_length} &
Length of the name inside the resource record &
Exfiltration techniques encode data into the name, which pushes its length
toward the protocol's upper limit. \\

\texttt{len} &
Total length of the domain name (FQDN) &
Same reasoning: an unusually long name is a classic tunneling signal. \\

\texttt{FQDN\_count} &
Number of components in the full domain name &
An unusual chain of subdomains can signal encoded data rather than a
legitimate naming hierarchy. \\

\texttt{subdomain\_length} &
Length of the subdomain portion &
Exfiltration typically encodes data in the subdomain rather than the
registered domain. \\

\texttt{subdomain} &
Whether a subdomain is present (derived indicator) &
Reinforces the signal already carried by \texttt{subdomain\_length}. \\

\texttt{labels} &
Number of labels (dot-separated segments) in the name &
An unusual number of segments is a structural hint of encoding. \\

\texttt{labels\_max} &
Length of the longest single label &
An abnormally long label rarely occurs in ordinary human naming. \\

\texttt{labels\_average} &
Average label length &
A complementary signal, less sensitive to a single extreme label than
\texttt{labels\_max}. \\

\texttt{numeric} &
Share of numeric characters in the name &
Base32/base64/hex-encoded data leaves an unusual proportion of digits
compared to a normal domain name. \\

\texttt{special} &
Share of special characters in the name &
Same reasoning as \texttt{numeric}: encoded data leaves a statistical trace
in the name's character composition. \\

\texttt{rr\_count} &
Number of resource records in the DNS response &
An unusual response volume can signal a channel being used to carry more
data per query than normal. \\

\texttt{a\_records} &
Number of A records in the response &
Some exfiltration or C2 techniques abuse A records to carry data or signal
status. \\

\texttt{ttl\_mean} &
Average TTL of the DNS records observed &
Infrastructure tied to malicious activity often uses very short TTLs
(fast-flux) or otherwise atypical TTLs compared to legitimate traffic. \\

\texttt{ttl\_variance} &
Variance of TTL across observed records &
High TTL variance can betray an unstable or deliberately shifting
infrastructure, typical of attack infrastructure. \\
\bottomrule
\end{tabular}
\caption{The 16 features NetScope uses to score a DNS query.}
\end{table}

These 16 columns are convenient during training because the dataset already
provides them. A real firewall or proxy, in production, only ever hands
over a domain name --- it does not pre-compute Shannon entropy or TTL
variance. \texttt{live\_features.py} closes this gap, treating the 16
features as two families: \textbf{10 name-only features}
(\texttt{entropy}, \texttt{len}, \texttt{FQDN\_count},
\texttt{subdomain\_length}, \texttt{subdomain}, \texttt{numeric},
\texttt{special}, \texttt{labels}, \texttt{labels\_max},
\texttt{labels\_average}), computable from the domain string alone, and
\textbf{6 DNS-response features} (\texttt{rr\_name\_entropy},
\texttt{rr\_name\_length}, \texttt{rr\_count}, \texttt{a\_records},
\texttt{ttl\_mean}, \texttt{ttl\_variance}), which require an actual DNS
resolution against a public resolver and fall back to 0 --- rather than
raising --- if \texttt{dnspython} is missing or the resolution fails.
The module's own documentation is explicit about the one feature it cannot
faithfully reconstruct: \texttt{ttl\_variance} is always 0 on this live
path, because a single DNS resolution only ever yields a single TTL value
--- a true variance would need several observations spread over time, which
this synchronous request/response flow does not have.

\subsection*{Feature Engineering and Selection}
The notebook does not simply fix the 16 features up front and move on. It
ranks them by importance (using the \texttt{feature\_importances\_} of a
trained Random Forest), then re-runs the comparison keeping only the top 8
to see how much performance is lost.

\begin{table}[h]
\centering
\small
\begin{tabular}{@{}lc@{}}
\toprule
\textbf{Feature} & \textbf{Importance (Random Forest)} \\
\midrule
\texttt{FQDN\_count}        & 0.174 \\
\texttt{subdomain\_length}  & 0.133 \\
\texttt{special}            & 0.126 \\
\texttt{numeric}            & 0.124 \\
\texttt{ttl\_mean}          & 0.121 \\
\texttt{labels}             & 0.115 \\
\texttt{labels\_average}    & 0.042 \\
\texttt{labels\_max}        & 0.029 \\
\texttt{rr\_count}          & 0.029 \\
\texttt{rr\_name\_entropy}  & 0.026 \\
\texttt{ttl\_variance}      & 0.025 \\
\texttt{rr\_name\_length}   & 0.022 \\
\texttt{subdomain}          & 0.015 \\
\texttt{len}                & 0.010 \\
\texttt{entropy}            & 0.009 \\
\texttt{a\_records}         & 0.000 \\
\bottomrule
\end{tabular}
\caption{Feature importance from the Random Forest trained in the notebook.}
\end{table}

Dropping to the top 8 features barely moves the result: F1 = 0.8309 with all
16 features against F1 = 0.8301 with only 8. Six features
(\texttt{FQDN\_count}, \texttt{subdomain\_length}, \texttt{special},
\texttt{numeric}, \texttt{ttl\_mean}, \texttt{labels}) already carry more
than 79\,\% of the total importance. The model that actually gets exported
--- and the one \texttt{scorer.py} loads in production --- keeps all 16
features rather than this reduced set, which is why \texttt{live\_features.py}
still has to reconstruct all 16 and not just the top 8.

\subsection*{Machine Learning Model}
Several model families were tested and compared. A first, simple pass
(Random Forest and XGBoost on a single train/test split) gives close
results (accuracy 0.835 for both). Before going further, the notebook
checks whether a trained model is even worth the trouble, by testing a
naive rule: flag a domain as suspicious if its entropy exceeds 3.5. That
rule reaches an accuracy of 0.582 but an F1 of only 0.0236 (precision
0.388, recall 0.012) --- it misses almost every real attack. On the same
test set, the trained models clear 0.83 F1, which is the honest answer to
whether the model --- and the entire production pipeline built around it
--- earns its complexity.

The notebook then compares five models with 5-fold stratified
cross-validation on F1 and ROC-AUC:

\begin{table}[h]
\centering
\small
\begin{tabular}{@{}lcccc@{}}
\toprule
\textbf{Model} & \textbf{F1 mean} & \textbf{F1 (std)} & \textbf{ROC-AUC mean} & \textbf{ROC-AUC (std)} \\
\midrule
LightGBM              & 0.8312 & 0.0017 & 0.8704 & 0.0012 \\
XGBoost                & 0.8309 & 0.0016 & 0.8707 & 0.0012 \\
Random Forest           & 0.8309 & 0.0017 & 0.8698 & 0.0011 \\
HistGradientBoosting    & 0.8308 & 0.0017 & 0.8707 & 0.0009 \\
Logistic Regression     & 0.8221 & 0.0015 & 0.8685 & 0.0021 \\
\bottomrule
\end{tabular}
\caption{Five-fold stratified cross-validation across five supervised models.}
\end{table}

Worth stating plainly: the four tree/boosting models are practically
indistinguishable --- the gap between the best and the worst of the four is
under 0.05 F1 points, smaller than their own standard deviation. Only
logistic regression, the simplest model, clearly trails behind. Picking
LightGBM owes more to a marginal edge and training/inference speed than to
any real performance gap over the other tree-based models. This choice also
has a direct production consequence: \texttt{NetScopeScorer} builds a
\texttt{shap.TreeExplainer}, which works cleanly for LightGBM (and the
other tree/boosting models tested) but would fail for logistic regression
--- one more reason the tree-based family, not just LightGBM specifically,
was the right lane to stay in.

LightGBM's hyperparameters were then tuned with \textbf{Optuna} (Bayesian
search, 50 trials) rather than an exhaustive grid --- a \texttt{GridSearchCV}
had already been used earlier in the notebook, but only for a preliminary
Random Forest tuning pass, and the notebook itself notes that grid search is
``expensive and coarse'' compared to a Bayesian search for the final tuning
step. Optuna's best hyperparameters: \texttt{n\_estimators=112},
\texttt{max\_depth=6}, \texttt{learning\_rate$\approx$0.144},
\texttt{num\_leaves=51}, \texttt{min\_child\_samples=23},
\texttt{subsample$\approx$0.946}, \texttt{colsample\_bytree$\approx$0.547},
\texttt{reg\_alpha$\approx$1.01}, \texttt{reg\_lambda$\approx$9.65e-8}, with
a cross-validated F1 of \textbf{0.8313}. The tuned model beat the base
LightGBM in the comparison, so it is the one exported (the notebook prints:
``Best model: tuned LightGBM (CV F1 = 0.8313)''). On the held-out test set,
its classification report reads:

\begin{table}[h]
\centering
\small
\begin{tabular}{@{}lcccc@{}}
\toprule
\textbf{Class} & \textbf{Precision} & \textbf{Recall} & \textbf{F1} & \textbf{Support} \\
\midrule
0 (Benign) & 1.00 & 0.72 & 0.84 & 12,018 \\
1 (Attack) & 0.72 & 1.00 & 0.83 & 8,535 \\
\midrule
Accuracy & \multicolumn{4}{c}{0.84} \\
\bottomrule
\end{tabular}
\caption{Classification report of the final tuned LightGBM model on the test set.}
\end{table}

Read honestly, this table means the model misses zero attacks on this test
set (recall = 1.00 on the Attack class) at the cost of a real false-positive
rate (precision = 0.72, so roughly 28\,\% of benign traffic gets flagged).
In a security context, trading precision for recall this way is generally
the right call --- missing a real exfiltration attempt is more costly than
chasing a false alarm, though this reading is not stated as a deliberate
design goal anywhere in the notebook itself and should be treated as an
observation made after the fact, not a documented objective. It also
foreshadows a production choice: with a model this willing to flag benign
traffic, the service cannot afford to act automatically on every
``suspect'' verdict --- exactly why the production code, not the notebook,
introduces a second, much higher bar before anything happens without a
human (see Risk Scoring and Decision, below).

The final model and its 16-feature list are exported with \texttt{joblib}
(\texttt{netscope\_model.pkl}, \texttt{netscope\_features.pkl}). This is the
handoff point between the two halves of NetScope: everything before this
line happened once, offline, on Kaggle; everything after it happens
continuously, in \texttt{src/netscope\_agent/}, every time a query is
scored.

\subsection*{Risk Scoring and Decision}
\texttt{NetScopeScorer.\_\_init\_\_} loads both exported files once at
service startup (\texttt{joblib.load}) and raises a clear
\texttt{FileNotFoundError} pointing to \texttt{models/README.md} if they are
missing, rather than failing obscurely later. \texttt{score\_domain()} then
does exactly what the notebook's final classification report measured:
\texttt{proba = model.predict\_proba(x)[0][1]}, turned into
\texttt{ml\_score = round(proba * 100, 1)} --- the single number every
other stage of the pipeline (thresholds, SHAP, the LLM, the dashboard) is
built around.

The production code, not the notebook, is where the actual decision logic
lives, and it uses \textbf{two separate thresholds} that should not be
confused with each other:

\begin{itemize}[nosep]
\item \textbf{50.0} --- the verdict threshold (\texttt{NetScopeScorer.threshold} in \texttt{scorer.py}): above it, a query is labeled \texttt{"suspect"} instead of \texttt{"normal"} for display purposes.
\item \textbf{80.0} --- \texttt{AUTO\_ACT\_THRESHOLD} in \texttt{auto\_scan.py}: above it, the system isolates the host and notifies CrisisBot \emph{automatically}, without waiting for a human to click anything. This same threshold and constant gate both the autonomous pull loop and the push webhook (\texttt{api.py} imports \texttt{AUTO\_ACT\_THRESHOLD} directly from \texttt{auto\_scan.py} rather than redefining it).
\end{itemize}

A third, purely cosmetic level exists only for the dashboard:
\texttt{risk\_level()} buckets the score into four labels (Low $<$30, Medium
$<$60, High $<$80, Critical $\geq$80), with no effect on any decision.

The three-term composite score explored in the training notebook (0.7 ×
ML score + 0.2 × traffic score + 0.1 × human score) \textbf{does not exist
in the production code}. The score used everywhere in the running service
is simply \texttt{ml\_score = round(proba × 100, 1)}. A user's simulated
human risk profile --- looked up deterministically from a fictional
directory, \texttt{data/users.json}, via a SHA-256 hash of the hostname
(\texttt{get\_user\_for\_host} in \texttt{actions.py}), standing in for a
real HR or phishing-simulation platform such as GoPhish or KnowBe4 ---
shapes the wording of the recommendation shown to the analyst (through
\texttt{human\_priority()}, which buckets the profile's own risk score into
``élevée'' $\geq$70, ``moyenne'' $\geq$40, or ``faible'' below that), not a
blended numeric score.

\subsection*{SHAP Explainability}
This is a case where the notebook and the production code disagree, and the
production code is the one that matters here. SHAP never appears in the
training notebook --- but the running service genuinely uses it.
\texttt{NetScopeScorer} builds a \texttt{shap.TreeExplainer(self.model)} at
startup, and every scored query goes through an \texttt{explain()} call
that returns the top 4 features driving that particular decision, each
tagged with a direction (increases or decreases risk). These get translated
into analyst-readable labels via a fixed dictionary in \texttt{scorer.py}
(\texttt{"rr\_name\_entropy"} $\rightarrow$ \emph{``Entropie du nom de
domaine élevée''}) and shown on the dashboard as ``Main factors''. The code
degrades gracefully (\texttt{self.explainer = None}, and \texttt{explain()}
then returns an empty list instead of raising) if \texttt{TreeExplainer}
fails to build --- relevant for a model like logistic regression, tested in
the model comparison above but not the one ultimately kept, which is
precisely why that comparison mattered for more than just picking the
highest F1.

\subsection*{Role of the LLM}
NetScope also calls an LLM (Groq, \texttt{llama-3.1-8b-instant}), but
strictly for writing, never for deciding --- the system prompt forbids it
explicitly: ``Ne jamais reclasser l'alerte (suspect/normal) : c'est déjà
décidé en amont'' (\emph{never reclassify the alert: that decision has
already been made upstream}). The LLM receives the score already computed,
the SHAP factors already computed, and the human risk profile already
looked up, and produces only a natural-language explanation and a
recommended action, returned as strict JSON with exactly two keys
(\texttt{explanation}, \texttt{recommendation}). With no Groq API key
configured, or if the call fails or returns something outside that shape, a
template-based fallback (\texttt{\_template\_fallback} in \texttt{llm.py})
builds the same kind of explanation from the same real data --- the same
SHAP factors, the same score, the same human-priority bucket --- never
inventing anything. This role --- explain, never decide --- lines up with
the principle already stated in Chapter~2, and puts NetScope alongside
DarkWatch's integrated mode and PhishHunter's deterministic path as places
where a language model narrates a decision it never made.

\subsection*{Continuous Monitoring and Memory}
NetScope does not use a learned memory either, only a lightweight
in-process cache, and like DarkWatch the exact mechanism differs by entry
point. \texttt{auto\_scan.py} keeps an in-memory set, \texttt{\_already\_acted},
of domains already isolated automatically during the current run, so the
same domain is not re-isolated on every subsequent draw of the sample; the
push webhook (\texttt{/webhook/traffic-event}) checks and updates the same
module-level set, so a domain isolated via one entry point is not
re-isolated via the other. This set --- exactly like DarkWatch's in-memory
deduplication set for its background auto-scan loop --- is lost on
restart: nothing here is persisted to disk, so a fresh process re-evaluates
every domain as if seeing it for the first time. On-demand scoring through
\texttt{/analyze} keeps no memory at all and never triggers automatic
isolation regardless of score, by design --- that endpoint only ever
returns a score for the dashboard to display.

\subsection*{Integration with CrisisBot}
A single shared function, \texttt{actions.record\_isolation}, logs the
incident (\texttt{data/incidents.json}, with \texttt{domain}, \texttt{user},
\texttt{ml\_score}, a timestamp, and a \texttt{source} tag) and notifies
CrisisBot, regardless of whether it was triggered by a human click
(\texttt{POST /isolate}), the autonomous loop, or the push webhook --- so
the incident only gets logged once no matter which of the three entry
points fired it. The notification itself (\texttt{crisisbot\_client.py})
follows the same best-effort pattern seen with DarkWatch and the other
detection agents: a plain \texttt{POST} to \texttt{CRISISBOT\_URL}
(\texttt{/api/webhook/incident}), a 3-second timeout, and \textbf{no
authentication of any kind on either side of the call} --- the exception is
caught, a message is printed to the console, and NetScope continues
operating normally, with no retry, exactly the same limitation already
flagged for DarkWatch rather than a NetScope-specific gap.

\subsection*{From Data to Display}
The NetScope dashboard (\texttt{static/netscope.html}) opens with a dark
``traffic card'' labeled as coming from a firewall/proxy webhook
(Fortinet, PaloAlto are named as examples), listing domains fetched from
\texttt{GET /sample-traffic?n=4} with their \texttt{ml\_score} shown as a
percentage; rows flagged \texttt{"suspect"} are highlighted in red with a
warning icon. The page always focuses on the first suspect result in that
sample: its domain, its NetScope risk score and colored level badge
(Faible/Moyen/Élevé/Critique), and its SHAP ``main factors'' list, each
prefixed with a \texttt{+} or \texttt{--} sign for the direction it pushed
the score. Below that, a profile card shows the correlated user's name,
department, and human risk score, followed by a one-sentence
correlation message whose wording depends on that user's priority bucket
(élevée / moyenne / faible) --- the same three-way bucket produced by
\texttt{human\_priority()} described above under Risk Scoring and
Decision. A single
``Isoler le poste'' button calls \texttt{POST /isolate} and, on success,
confirms the exact timestamp and states explicitly that the incident was
logged to \texttt{data/incidents.json} --- the same shared function
described above, this time with \texttt{source="manual"}.

\subsection*{Limitations}
All five supervised models compared share the same structural limit: they
only recognize attacks that statistically resemble what is already present
in CIC-Bell-DNSExF2021. A genuinely new DNS exfiltration technique, never
represented in that dataset, could slip through undetected --- a limit of
the offline training data that no amount of production engineering around
it can change. \texttt{live\_features.py} documents an honest limitation of
its own on the production side: \texttt{ttl\_variance} is always 0 on the
live path, since a single DNS resolution only yields a single TTL, and its
reconstruction of the DNS response features is a reasonable approximation,
not a certified match to the original dataset authors' exact methodology
--- a gap between the offline features the model learned from and the
online features it is actually served in production, even though both are
computed from the same 16 definitions. Neither the CrisisBot call nor the
dashboard itself is authenticated, the same limitation already noted for
DarkWatch. The project's own \texttt{README.md} still lists verifying the
absence of data leakage in the features among its next steps --- something
this audit did not re-verify and that remains an open question.

\subsection*{Conclusion}
NetScope pairs a conventional but carefully built machine learning pipeline
--- five model families compared under cross-validation, Bayesian
hyperparameter tuning, and an explicit sanity check against a naive rule
--- with a production layer that closes the gap between a research dataset
and a real domain name (\texttt{live\_features.py}), adds real SHAP
explainability, and lets an LLM write but never decide. The two thresholds
that actually drive behavior, 50 for the suspect/normal verdict and 80 for
automatic isolation, live entirely in the service code, not the training
notebook, and neither matches a single ``80/100'' figure loosely used
elsewhere. As a supervised classifier, NetScope remains bound by what its
training data has already seen, a limitation the project's own code and
documentation acknowledge rather than hide. The next section turns to
PhishHunter, the only other agent in \ProjectName{} built around a trained
model, and the one where --- unlike NetScope's LLM, confined to writing ---
a language model is deliberately allowed a second, genuinely agentic
decision path of its own.
